\chapter{Modello dati}

Tutto vive in \texttt{webui/backend/models.py}. Le tabelle si
raggruppano in cinque famiglie tematiche.

\section{Famiglie di tabelle}

\subsection*{Anagrafiche}
\texttt{subjects}, \texttt{teachers}, \texttt{school\_classes},
\texttt{classrooms}, \texttt{curricula}, \texttt{students},
\texttt{study\_groups}.

\subsection*{Tag (M-N)}
\texttt{classroom\_tags} + \texttt{classroom\_tag\_assignments},
\texttt{student\_tags} + \texttt{student\_tag\_assignments}. Nome
unique lower-cased, descrizione opzionale, cancellazione cascata.

\subsection*{Assegnazione}
\texttt{class\_subjects}, \texttt{curriculum\_subject\_hours},
\texttt{teacher\_subjects}, \texttt{subject\_group\_weights},
\texttt{assignments}.

\subsection*{Disponibilit\`a / vincoli}
\texttt{teacher\_unavailability}, \texttt{class\_unavailability},
\texttt{classroom\_unavailability},
\texttt{logical\_unavailabilities},
\texttt{curriculum\_logical\_constraints},
\texttt{classroom\_subject\_preferences},
\texttt{teacher\_classroom\_preferences},
\texttt{coteaching\_rules}.

\subsection*{Soluzioni e scheduling}
\texttt{solutions}, \texttt{lessons}, \texttt{day\_counts}.

\subsection*{Coverage e Run}
\texttt{absences}, \texttt{substitute\_assignments},
\texttt{runs}, \texttt{run\_logs}, \texttt{app\_state}.

\subsection*{Viste salvate}
\texttt{saved\_views}: query DSL + sort multi-livello salvati per
una entit\`a lista, esposti dal dropdown "Viste salvate" del
componente \emph{SortableQueryableList}. UNIQUE \texttt{(entity,
name)}.

\section{Diagramma relazioni}

\begin{figure}[h]
\centering
\begin{tikzpicture}[node distance=0.5cm and 0.4cm,
  rect/.style={draw, rounded corners=2pt, thick, fill=cBox!20,
               minimum width=2.4cm, minimum height=0.6cm,
               align=center, font=\footnotesize}]
  \node[rect] (sub) {subjects};
  \node[rect, right=of sub] (ts) {teacher\_subjects};
  \node[rect, right=of ts] (te) {teachers};
  \node[rect, below=of te] (asn) {assignments};
  \node[rect, left=of asn] (sc) {school\_classes};
  \node[rect, left=of sc] (cur) {curricula};
  \node[rect, below=of asn] (sol) {solutions};
  \node[rect, below=of sol] (les) {lessons};
  \node[rect, below=of cur] (cs) {class\_subjects};
  \node[rect, below=of sc] (cu) {class\_unavail.};
  \node[rect, right=of les] (rm) {classrooms};
  \node[rect, above=of rm] (cp) {classroom\_subj\_pref};
  \draw (sub) -- (ts) -- (te);
  \draw (te) -- (asn);
  \draw (sc) -- (asn);
  \draw (sc) -- (cs);
  \draw (sc) -- (cu);
  \draw (cur) -- (sc);
  \draw (asn) -- (sol);
  \draw (sol) -- (les);
  \draw (rm) -- (cp);
  \draw (sub) to[out=270,in=180,looseness=1.4] (cp);
\end{tikzpicture}
\caption{Relazioni rilevanti del modello dati (semplificato).}
\label{fig:datamodel}
\end{figure}

\section{Migrazioni idempotenti}

Pattern (in \texttt{db.py}):

\begin{lstlisting}[style=py]
if insp.has_table("school_classes") \
        and not has_column("school_classes", "curriculum_id"):
    conn.execute(text(
        "ALTER TABLE school_classes ADD COLUMN curriculum_id INTEGER"
    ))
\end{lstlisting}

Ogni ALTER \`e guardato da una check su \texttt{PRAGMA table\_info}.
Quando aggiungi una colonna nuova, segui lo stesso pattern: il DB
preesistente sopravvive senza interventi manuali.

\section{Pickle dell'engine}

\texttt{engine\_io.py} converte fra DB e tre forme pickle:

\begin{lstlisting}[style=py]
school = {
  'profile':   str,
  'classes':   [{name, year, section, curriculum, subjects:{subj:ore}}],
  'teachers':  [{name, group, max_hours, free_day,
                 weights:{subj:int}}],
  'cconcorsopersubject': {subj: {group: weight}},
  'curriculum_scores':   {curriculum: int},
  'curricula':           [...],
  'students':            [...],
  'groups':              [...],
}
profs = {
  teacher_name: {
    'classi':  {class_name: {subject: {'ore': N}}},
    'glibero': [d1, d2, d3]
  }
}
solution = {(prof, class, subj, day, hour): 0|1}
\end{lstlisting}

% =========== Cap 4: Vincoli ===========
